% ============================================================
%  preamble.tex — The Secret Life of the Internet
%  Shared packages, colours, fonts, and reusable visual
%  callout macros for a kid-friendly, visual-first LaTeX book.
% ============================================================


% ============================================================
%  DIGITAL EDITION SETTINGS
%  - Page: 8.5×11 (comfortable on tablets and screens)
%  - Narrow margins for maximum readable area
%  - sRGB colour output (screen standard)
%  - Active hyperlinks with PDF bookmarks
%  - Tagged PDF structure for accessibility
%  - Increased line spacing for screen readability
% ============================================================
% --- Core packages ---
\usepackage[T1]{fontenc}
\usepackage[utf8]{inputenc}
\usepackage[sfdefault]{sourcesanspro}
\usepackage[
  paperwidth=8.5in,          % Wider page for screen/tablet reading
  paperheight=11in,
  margin=0.75in,             % Narrow margins — more content area on screen
  top=0.9in,
  bottom=0.8in
]{geometry}
\usepackage{xcolor}
\usepackage{graphicx}
\usepackage{tcolorbox}
\tcbuselibrary{skins,breakable}
\usepackage{tikz}
\usetikzlibrary{arrows.meta,shapes,shapes.symbols,positioning,decorations.pathmorphing,calc,shadows,fit,backgrounds}
\usepackage{fontawesome5}    % icon glyphs
\usepackage{mdframed}
\usepackage{parskip}
\usepackage{setspace}
\setstretch{1.15}            % Slightly increased line spacing for screen reading
\usepackage{enumitem}
\usepackage{array}
\usepackage{tabularx}
\usepackage{ragged2e}
\usepackage{needspace}
\usepackage{etoolbox}
% microtype: protrusion only — expansion disabled to avoid MiKTeX
% "auto expansion is only possible with scalable fonts" pdfTeX error.
\usepackage[protrusion=true,expansion=false]{microtype}
\usepackage{fancyhdr}

% --- Book-wide colour palette ---
% Colours must be defined before \hypersetup uses them.
\definecolor{InternetBlue}{HTML}{2196F3}
\definecolor{WarmOrange}{HTML}{FF9800}
\definecolor{LeafGreen}{HTML}{4CAF50}
\definecolor{SunYellow}{HTML}{FFC107}
\definecolor{CoralRed}{HTML}{F44336}
\definecolor{SoftPurple}{HTML}{9C27B0}
\definecolor{SkyBlue}{HTML}{87CEEB}
\definecolor{LightGray}{HTML}{F5F5F5}
\definecolor{DarkText}{HTML}{212121}

\usepackage{hyperref}
\hypersetup{
  colorlinks=true,
  linkcolor=InternetBlue,
  urlcolor=InternetBlue,
  citecolor=InternetBlue,
  bookmarks=true,            % Generate PDF bookmarks (clickable TOC)
  bookmarksnumbered=true,    % Chapter numbers in bookmarks
  bookmarksopen=true,        % Expand bookmarks by default
  pdfstartview=FitH,         % Fit page width on open
  pdftitle={The Secret Life of the Internet},
  pdfauthor={Fansci Solutions},
  pdfsubject={Technology Explorer Series — Digital Edition}
}

% --- Accessibility: tagged PDF and alt-text support ---
\usepackage{accsupp}
\usepackage{accessibility}   % Tagged PDF output for screen readers
\newcommand{\alttext}[2]{%
  \BeginAccSupp{method=escape,ActualText={#1}}%
  #2%
  \EndAccSupp{}%
}

% --- Index support ---
\usepackage{makeidx}
\makeindex

% --- Chapter title styling ---
\usepackage{titlesec}
\titleformat{\chapter}[display]
  {\normalfont\huge\bfseries\color{InternetBlue}}
  {\chaptertitlename\ \thechapter}{12pt}{\Huge}
\titlespacing*{\chapter}{0pt}{14pt}{16pt}
\titleformat{\section}
  {\normalfont\LARGE\bfseries\color{InternetBlue}}{\thesection}{0.6em}{}
\titleformat{\subsection}
  {\normalfont\Large\bfseries\color{WarmOrange}}{\thesubsection}{0.6em}{}

\AtBeginDocument{\color{DarkText}}

% --- Header / footer style ---
% Keep headheight large enough for the coloured running header text.
\setlength{\headheight}{14pt}
\setlength{\emergencystretch}{2em}
\pagestyle{fancy}
\fancyhf{}
\fancyhead[L]{\nouppercase{\small\bfseries\color{InternetBlue}\leftmark}}
\fancyfoot[C]{\color{InternetBlue}\thepage}
\renewcommand{\headrulewidth}{0.6pt}
\renewcommand{\footrulewidth}{0pt}
\fancypagestyle{plain}{%
  \fancyhf{}
  \fancyfoot[C]{\color{InternetBlue}\thepage}
  \renewcommand{\headrulewidth}{0pt}
  \renewcommand{\footrulewidth}{0pt}
}
\makeatletter
\renewcommand{\cleardoublepage}{%
  \clearpage
  \if@twoside
    \ifodd\c@page\else
      \hbox{}%
      \thispagestyle{empty}%
      \newpage
      \if@twocolumn\hbox{}\newpage\fi
    \fi
  \fi
}
\makeatother
\preto{\section}{\needspace{8\baselineskip}}
\preto{\subsection}{\needspace{6\baselineskip}}

% ============================================================
%  REUSABLE CALLOUT BOXES
% ============================================================

% --- Big Idea box (blue) ---
\newtcolorbox{bigidea}[1][]{
  enhanced,
  before upper=\RaggedRight,
  colback=InternetBlue!10,
  colframe=InternetBlue,
  fonttitle=\bfseries\large,
  title={\faLightbulb\quad Big Idea},
  arc=6pt,
  boxrule=1.5pt,
  left=8pt, right=8pt, top=6pt, bottom=6pt,
  #1
}

% --- Picture This box (orange) ---
\newtcolorbox{picturethis}[1][]{
  enhanced,
  before upper=\RaggedRight,
  colback=WarmOrange!10,
  colframe=WarmOrange,
  fonttitle=\bfseries\large,
  title={\faPaintBrush\quad Picture This},
  arc=6pt,
  boxrule=1.5pt,
  left=8pt, right=8pt, top=6pt, bottom=6pt,
  #1
}

% --- Try It box (green) ---
\newtcolorbox{tryit}[1][]{
  enhanced,
  before upper=\RaggedRight,
  colback=LeafGreen!10,
  colframe=LeafGreen,
  fonttitle=\bfseries\large,
  title={\faHandPointRight\quad Try It},
  arc=6pt,
  boxrule=1.5pt,
  left=8pt, right=8pt, top=6pt, bottom=6pt,
  #1
}

% --- Quick Map box (yellow) ---
\newtcolorbox{quickmap}[1][]{
  enhanced,
  before upper=\RaggedRight,
  colback=SunYellow!15,
  colframe=SunYellow!80!black,
  fonttitle=\bfseries\large,
  title={Quick Map},
  arc=6pt,
  boxrule=1.5pt,
  left=8pt, right=8pt, top=6pt, bottom=6pt,
  #1
}

% --- New Words box (purple) ---
\newtcolorbox{newwords}[1][]{
  enhanced,
  before upper=\RaggedRight,
  colback=SoftPurple!8,
  colframe=SoftPurple,
  fonttitle=\bfseries\large,
  title={\faBook\quad New Words},
  arc=6pt,
  boxrule=1.5pt,
  left=8pt, right=8pt, top=6pt, bottom=6pt,
  #1
}

% --- Fun Fact box (coral red) ---
\newtcolorbox{funfact}[1][]{
  enhanced,
  before upper=\RaggedRight,
  colback=CoralRed!8,
  colframe=CoralRed,
  fonttitle=\bfseries\large,
  title={\faStar\quad Fun Fact},
  arc=6pt,
  boxrule=1.5pt,
  left=8pt, right=8pt, top=6pt, bottom=6pt,
  #1
}


% --- Digital: responsive diagram scaling for screen ---
\newcommand{\screendiagram}[2][0.95]{%
  \begin{center}\resizebox{#1\linewidth}{!}{#2}\end{center}%
}
% Tip: wrap TikZ in \screendiagram for auto-scaling on different screen sizes.
% ============================================================
%  DIAGRAM WRAPPER (with accessibility alt-text)
%  Usage: \begin{diagram}[Optional caption / alt-text]
%           ... TikZ or other content ...
%         \end{diagram}
%
%  The optional argument serves as both the visible caption
%  AND the accessible ActualText for screen readers.
%  \makeatletter / \makeatother are required because the
%  environment uses internal @ macros (\@diagramcaption,
%  \@empty).
% ============================================================
\makeatletter
\newenvironment{diagram}[1][]{%
  \def\@diagramcaption{#1}%
  \begin{figure}[htbp]
    \centering
    \BeginAccSupp{method=escape,ActualText={#1}}%
    \begin{tcolorbox}[
      enhanced,
      colback=SkyBlue!8,
      colframe=SkyBlue!50!black,
      boxrule=0.8pt,
      arc=6pt,
      left=8pt, right=8pt, top=8pt, bottom=8pt
    ]
    \centering
}{%
    \end{tcolorbox}
    \EndAccSupp{}%
    \ifx\@diagramcaption\@empty\else
      \vspace{4pt}{\small\itshape\@diagramcaption}%
    \fi
  \end{figure}%
}
\makeatother

% ============================================================
%  CHARACTER CALLOUT HELPER
%  Usage: \character{Dot}{Let's follow me through the internet!}
% ============================================================
\newcommand{\character}[2]{%
  \begin{tcolorbox}[
    enhanced,
    colback=SkyBlue!15,
    colframe=SkyBlue!60!black,
    arc=10pt,
    boxrule=1pt,
    left=10pt, right=10pt, top=6pt, bottom=6pt,
    title={\textbf{\Large #1}},
    fonttitle=\bfseries\color{DarkText},
    before upper=\RaggedRight
  ]
  \itshape ``#2''
  \end{tcolorbox}%
}

% ============================================================
%  CHAPTER OPENING HELPER
%  Usage: \chapteropening{Big question the chapter answers.}
% ============================================================
\newcommand{\chapteropening}[1]{%
  \begin{tcolorbox}[
    enhanced,
    colback=LightGray,
    colframe=InternetBlue!40,
    arc=4pt,
    boxrule=0.8pt,
    left=10pt, right=10pt, top=8pt, bottom=8pt,
    before upper=\RaggedRight
  ]
  \large\itshape #1
  \end{tcolorbox}%
  \vspace{6pt}%
}

% --- Chapter badge helper ---
\newcommand{\chapterbadge}[1]{%
  \begin{center}
    \begin{tcolorbox}[enhanced,colback=SunYellow!20,colframe=SunYellow!70!black,arc=10pt,boxrule=1pt,left=10pt,right=10pt,top=8pt,bottom=8pt]
      \centering\Large\bfseries\color{InternetBlue}#1
    \end{tcolorbox}
  \end{center}
}

% --- Small utility: bold coloured term inline (also adds index entry) ---
\newcommand{\term}[1]{\textbf{\color{InternetBlue}#1}\index{#1}}

% --- Glossary term helper: bold coloured term without an index entry ---
% Use \glossterm instead of \term inside the Glossary so that
% glossary pages do not flood the back-of-book index with duplicates.
\newcommand{\glossterm}[1]{\textbf{\color{InternetBlue}#1}}

% --- Horizontal rule between sections ---
\newcommand{\sectionrule}{\vspace{4pt}\noindent\textcolor{SkyBlue!60}{\rule{\linewidth}{0.6pt}}\vspace{4pt}}

% --- Table helpers ---
\newcolumntype{Y}{>{\RaggedRight\arraybackslash}X}
\newcolumntype{C}[1]{>{\centering\arraybackslash}p{#1}}

% --- Starred chapter helper with clean running marks ---
\newcommand{\namedchapter}[1]{%
  \chapter*{#1}%
  \markboth{#1}{}%
  \thispagestyle{plain}%
}
